\documentclass[10pt,a4paper]{amsart}
\usepackage[margin=19mm]{geometry}
\usepackage{amsmath,amssymb,mathtools}
\usepackage[scr=rsfs,cal=euler]{mathalfa}
\usepackage{xfrac}
\usepackage[unicode,hidelinks,bookmarks=false]{hyperref}
\newcommand{\ie}{i.e.\ }
\newcommand{\cf}{cf.\ }
\newcommand{\resp}{respectively\ }
\newcommand{\Subsection}{\subsection}
\providecommand{\nobreakdash}{\nobreak\hbox{-}}
\setlength{\emergencystretch}{2em}
\multlinegap0pt
\usepackage{amssymb}
\usepackage{dsfont}
\usepackage{mathtools}
\newtheorem{lemma}{Lemma}
\newcommand{\beq}{\begin{equation}}
\newcommand{\eeq}{\end{equation}}
\title{Convergence, Sticking and Escape: Stochastic Dynamics Near Critical Points in SGD}
\author{Dmitry Dudukalov, Artem Logachov, Vladimir Lotov, Timofei Prasolov, Evgeny Prokopenko, Anton Tarasenko}
\date{Source: arXiv:2505.18535v2 (CC BY 4.0)}
\begin{document}
\maketitle

\noindent\emph{Source and licence.} Convergence, Sticking and Escape: Stochastic Dynamics Near Critical Points in SGD. Version arXiv:2505.18535v2 (CC BY 4.0), distributed by arXiv under the Creative Commons Attribution 4.0 International licence, http://creativecommons.org/licenses/by/4.0/, as recorded in the arXiv licence metadata for this exact version and shown on the abstract page https://arxiv.org/abs/2505.18535v2. Full licence notice: \texttt{licenses/arxiv-2505.18535-CC-BY-4.0.txt}.\par\smallskip

\noindent\textit{Editorial scope.} Counted unit: the article's lemma l.1 (source0.tex lines 1340--1346). The article's complete original proof of it is delivered with it (source0.tex lines 1348--1408, one \texttt{proof} environment of 61 lines). No mathematics is written, repaired or re-derived: the statement and the proof are the article's own lines, in the article's own notation, and no retained line is substituted or re-flowed.  No further source-local context is needed: every cross-reference of every retained line resolves inside the two delivered spans.  Nothing was omitted from any retained span, so the package carries no omission record.  No retained line contains a citation, so the wrapper carries no bibliography and no bibliography entry.  No retained line contains a figure environment, an image-inclusion command or drawing code, so no image is delivered and the package declares no asset dependency.  Editorial formatting only, in addition: an AMS \texttt{amsart} preamble that restates the article's own declarations and macros used by the retained lines, namely the environments \texttt{lemma} on separate counters, together with the standard packages those lines need.  The retained environments print in this document as Lemma 1 in order of appearance, whereas the article prints the same items with the numbering of its own full document.  The complete native source of the article as distributed in the arXiv e-print for this exact version is bundled unchanged as \texttt{sources/178700/source0.tex}.
\begin{lemma} \label{l.1}
For $0<\delta<M_r-m$, $x\in[m,M_r-\delta)$,
$0<\varepsilon<\frac{\delta}{10 C_{\max}}$ the following equality holds:
\beq\label{20.11.5}
A_{\frac{\delta}{5}}\left(\varepsilon\right)\cap B_{x,\delta}\left(\varepsilon\right)=A_{\frac{\delta}{5}}\left(\varepsilon\right).
\eeq
\end{lemma}

\begin{proof}
If $A_{\frac{\delta}{5}}\left(\varepsilon\right)=\varnothing$ then it is easy to see that (\ref{20.11.5})
holds. Assume that $A_{\frac{\delta}{5}}\left(\varepsilon\right)\neq\varnothing$.
Let us define a sequence of stopping times:
$$
\tau_0:=0, \ \ \
\tau_{2v-1}:=\inf\{\tau_{2v-2}< k \leqslant n_\varepsilon:x_k<m\}\wedge \lfloor n_\varepsilon \rfloor,
$$
$$
\tau_{2v}:=\inf\{\tau_{2v-1}< k \leqslant n_\varepsilon:x_k\geqslant m\}\wedge \lfloor n_\varepsilon \rfloor.
$$
Here $1\leqslant v \leqslant v_{\max}$, $v_{\max}:=\min\{v\in \mathbb{N}:\tau_{2v}= \lfloor n_\varepsilon \rfloor\}$, $\inf\{\varnothing\}=\infty$.

Let us point out that on event $A_{\frac{\delta}{5}}\left(\varepsilon\right)$ the following equality holds:
\beq \label{21.10.1}
x_k^{\varepsilon}<x+\delta 
\eeq
for $k=\tau_0$ and $\tau_{2v-1}\leqslant k < \tau_{2v}$, $1\leqslant v \leqslant v_{\max}$.


It is easy to that on event $A_{\frac{\delta}{5}}\left(\varepsilon\right)$ the following holds:
\beq \label{21.10.2}
\varepsilon\left|\sum\limits_{j=l_1}^{l_2}\xi_j\right|<\frac{2\delta}{5}
\eeq
for every $1\leqslant l_1\leqslant l_2 \leqslant n_\varepsilon$. Hence for $0<\varepsilon<\frac{\delta}{10C_{\max}}$,
$1\leqslant v < v_{\max}$ inequality
\beq \label{21.10.3}
x_{\tau_{2v}}^{\varepsilon}\leqslant m+ \varepsilon C_{\max}+\varepsilon\xi_{\tau_{2v}}
\leqslant m+ \varepsilon C_{\max}+\frac{2\delta}{5}<m+\frac{\delta}{2}
\leqslant x+\frac{\delta}{2}
\eeq
holds.

Now we use proof by induction. Consider the case when
 $\tau_{2v-2}<k< \tau_{2v-1}$, $1\leqslant v < v_{\max}$. Since $f'(x_{\tau_{2v-2}}^{\varepsilon})>0$ it follows from definition of stopping time $\tau_0$
and formulae (\ref{21.10.2}), (\ref{21.10.3})  that on event $A_{\frac{\delta}{5}}\left(\varepsilon\right)$ the following inequality holds for $k=\tau_{2v-2}+1$, $1\leqslant v < v_{\max}$:
\beq \label{21.10.5}
x_k^{\varepsilon}=x_{\tau_{2v-2}}^{\varepsilon}-\varepsilon f'(x_{\tau_{2v-2}}^{\varepsilon})+\varepsilon\xi_{\tau_{2v-2}+1}
\leqslant x+\frac{\delta}{2}+\frac{2\delta}{5}<x+\delta.
\eeq

We fix $1\leqslant v < v_{\max}$.
We assume that for some $\tau_{2v-2}+1<l<\tau_{2v-1}$ inequality (\ref{21.10.5}) holds for every $\tau_{2v-2}+1\leqslant k\leqslant l$. Let us show that it holds for $k=l+1$. It follows from induction hypothesis that on event $A_{\frac{\delta}{5}}\left(\varepsilon\right)$ we have
\beq \label{21.10.7}
f'(x_k^{\varepsilon})>0
\eeq
for every $\tau_{2v-2}+1\leqslant k\leqslant l$.

From inequalities (\ref{21.10.2}), (\ref{21.10.3}), (\ref{21.10.7}) on event $A_{\frac{\delta}{5}}\left(\varepsilon\right)$ we have
$$
x_{l+1}^{\varepsilon}=x_{\tau_{2v-2}}^{\varepsilon}-\varepsilon \sum\limits_{j=\tau_{2v-2}}^{l}f'(x_j^{\varepsilon})
+\varepsilon\sum\limits_{j=\tau_{2v-2}+1}^{l+1}\xi_j
\leqslant x_{\tau_{2v-2}}^{\varepsilon}+\frac{2\delta}{5}<x+\delta.
$$
Hence we proved that induction hypothesis is true, i.e., on event $A_{\frac{\delta}{5}}\left(\varepsilon\right)$ inequality (\ref{21.10.5})
holds for every $\tau_{2v-2}+1\leqslant k< \tau_{2v-1}$ and every $1\leqslant v < v_{\max}$.

Lemma \ref{l.1} follows from inequalities (\ref{21.10.1}), (\ref{21.10.3})
and the fact that formula (\ref{21.10.5})
holds for every $\tau_{2v-2}+1\leqslant k< \tau_{2v-1}$, $1\leqslant v < v_{\max}$.
\end{proof}

\end{document}
